We have presented evidence that internal computational routing structures—as reflected in attention patterns—project to observable signals after dimensionality reduction. Specifically:

\begin{enumerate}
    \item Pathway metrics extracted from attention distributions monotonically relate to attention sparsity ($r = 0.68$), validating our metric extraction approach.
    
    \item EEG-projected signals retain moderate predictive power for sparsity ($r = 0.42$), demonstrating that information survives linear projection with noise injection.
    
    \item Pathway knowledge meaningfully transfers to EEG predictions (TSR = 0.62), suggesting a causal relationship between routing structure and population signals.
\end{enumerate}

These findings support the hypothesis that computational pathways project to non-invasive population-level recordings in a predictable way.

\subsection{Broader Significance}

From a neuroscience perspective, our work contributes to the interpretability of EEG. Rather than viewing EEG as direct measurements of stimulus-locked neural responses, we propose EEG encodes the organization of information routing—a potentially deeper and more stable marker of neural computation.

From an AI perspective, we demonstrate that mechanistic understanding of neural networks (via attention analysis) can predict behavior of data after realistic measurement transformations, bridging interpretability and predictive modeling.

\subsection{Final Remarks}

This work is a proof of concept in controlled conditions. Extending these findings to real neural data and biological relevance remains a frontier of research. We hope this framework encourages further investigation into the relationship between internal routing and population signals, with implications for both neuroscience and AI interpretability.
